%
% $Log: studguide.tex,v $
% Revision 1.5  2002/10/21 15:26:44  mjm
% to ohio, revised
%
% Revision 1.4  2001/08/15 18:07:03  mjm
% APPM preprint 466
%
% Revision 1.3  1999/08/09 19:55:27  mjm
% beta release
%
% Revision 1.2  1999/04/06 19:22:02  mjm
% Ellens comments
%
%

%%%%%%%%%%%%%%%%%%%%%%%%%%%%%%%%%%%%%%%%%%%%%%%%%%%%%%%%%%%%%%%%%%%%%%%%%
\handout{Student's Guide to Good Problems}

Many students believe that writing has nothing to
do with mathematics. {\bf This is false.}
In any field, it is important to be able to communicate your ideas and
results to others. Exactly how you communicate them varies from field to
field, but in nearly all fields it is through writing. Mathematical
writing, however, has its own particular style. The emphasis is on clarity
and precision, and not on the clever turn of phrase.

The ``Good Problems'' program is designed to teach you how to write
about mathematics coherently. It provides you a set of specific
writing skills and gives you enough practice through regular
assignments so that writing well will become a habit. It is also
designed to be as painless as possible.

Throughout the term, some homework problems will be designated 
``good problems''. You  need to write the solution to this problem
carefully, in good ``presentation'' format. It will be graded partly for
correctness, but with emphasis on presentation.

During the term you will receive six handouts, each addressing a
particular writing skill. The good problem is graded only on those
skills already covered, with greater emphasis on the newest skill. By
the end of the term you will have acquired the basic skills of
mathematical writing. The handouts are:
\begin{itemize}
\item Laying out the Problem.
\item Flow.
\item Mathematical Symbols.
\item Logical Connectives.
\item Graphs.
\item Introductions and Conclusions.
\end{itemize}
Each handout consists of a description of a particular writing skill,
motivation for the importance of this skill,
whatever explicit rules can be given,
and examples of good and bad presentation emphasizing that skill.
We use a consistent format for the examples.

\smallskip
\begin{badexample}
An example of poor writing looks like this.
\end{badexample}

\smallskip
\begin{goodexample}
An example of good writing looks like this.\\
\fbox{\rm A comment within an example looks like this.}
\end{goodexample}

\smallskip
If you already have these basic writing skills, then it should only take you
ten extra minutes to write up a good problem. If you do not have
these skills, then you will have to spend extra effort in the short term to
acquire them. In the long term, however, these skills will save you much
time and effort. Besides making future writing assignments easier, these
skills will help you with problem-solving (by encouraging organization and
logic) and with reading mathematics. These skills should be carried into
future mathematics courses, and become a useful life skill.


